\documentclass[10pt,a4paper]{amsart}
\usepackage{amsmath,amssymb,amsthm,graphicx}
\usepackage{verbatim}
\newenvironment{reference}{\comment}{\endcomment}
%\newenvironment{reference}{}{}
\newenvironment{slogan}{\comment}{\endcomment}
\newenvironment{history}{\comment}{\endcomment}

% For commutative diagrams we use Xy-pic
\usepackage[all]{xy}

% We use 2cell for 2-commutative diagrams.
\xyoption{2cell}
% Theorem environments.
%
\theoremstyle{plain}
\newtheorem{theorem}[subsection]{Theorem}
\newtheorem{proposition}[subsection]{Proposition}
\newtheorem{lemma}[subsection]{Lemma}

\theoremstyle{definition}
\newtheorem{definition}[subsection]{Definition}
\newtheorem{example}[subsection]{Example}
\newtheorem{exercise}[subsection]{Exercise}
\newtheorem{situation}[subsection]{Situation}

\theoremstyle{remark}
\newtheorem{remark}[subsection]{Remark}
\newtheorem{remarks}[subsection]{Remarks}

\numberwithin{equation}{subsection}

% Macros
%
\def\lim{\mathop{\mathrm{lim}}\nolimits}
\def\colim{\mathop{\mathrm{colim}}\nolimits}
\def\Spec{\mathop{\mathrm{Spec}}}
\def\Hom{\mathop{\mathrm{Hom}}\nolimits}
\def\Ext{\mathop{\mathrm{Ext}}\nolimits}
\def\SheafHom{\mathop{\mathcal{H}\!\mathit{om}}\nolimits}
\def\SheafExt{\mathop{\mathcal{E}\!\mathit{xt}}\nolimits}
\def\Sch{\mathit{Sch}}
\def\Mor{\mathop{\mathrm{Mor}}\nolimits}
\def\Ob{\mathop{\mathrm{Ob}}\nolimits}
\def\Sh{\mathop{\mathit{Sh}}\nolimits}
\def\NL{\mathop{N\!L}\nolimits}
\def\CH{\mathop{\mathrm{CH}}\nolimits}
\def\proetale{{pro\text{-}\acute{e}tale}}
\def\etale{{\acute{e}tale}}
\def\QCoh{\mathit{QCoh}}
\def\Ker{\mathop{\mathrm{Ker}}}
\def\Im{\mathop{\mathrm{Im}}}
\def\Coker{\mathop{\mathrm{Coker}}}
\def\Coim{\mathop{\mathrm{Coim}}}

% Boxtimes
%
\DeclareMathSymbol{\boxtimes}{\mathbin}{AMSa}{"02}

%
% Macros for moduli stacks/spaces
%
\def\QCohstack{\mathcal{QC}\!\mathit{oh}}
\def\Cohstack{\mathcal{C}\!\mathit{oh}}
\def\Spacesstack{\mathcal{S}\!\mathit{paces}}
\def\Quotfunctor{\mathrm{Quot}}
\def\Hilbfunctor{\mathrm{Hilb}}
\def\Curvesstack{\mathcal{C}\!\mathit{urves}}
\def\Polarizedstack{\mathcal{P}\!\mathit{olarized}}
\def\Complexesstack{\mathcal{C}\!\mathit{omplexes}}
% \Pic is the operator that assigns to X its picard group, usage \Pic(X)
% \Picardstack_{X/B} denotes the Picard stack of X over B
% \Picardfunctor_{X/B} denotes the Picard functor of X over B
\def\Pic{\mathop{\mathrm{Pic}}\nolimits}
\def\Picardstack{\mathcal{P}\!\mathit{ic}}
\def\Picardfunctor{\mathrm{Pic}}
\def\Deformationcategory{\mathcal{D}\!\mathit{ef}}
\usepackage{cleveref}
\crefname{originallemma}{original lemma}{original lemmas}
\Crefname{originallemma}{Original lemma}{Original lemmas}
\makeatletter
\newlabel{lemma-gabriel-popescu-left-adjoint}{{0F5S}{source}{}{sourcepointer-1}{}}
\newlabel{lemma-gabriel-popescu-left-adjoint@cref}{{[originallemma][0][]0F5S}{[1][0][]source}{}{}{}}
\makeatother
\title{The left adjoint of the Gabriel-Popescu functor preserves injective maps}
\author{The Stacks Project Authors}
\date{stacks-project:injectives.tex:local-snapshot-20261012; GNU FDL 1.2 or later}
\begin{document}
\maketitle
\noindent\textit{Editorial note.} All delivered mathematical text is the author's own, copied from the pinned Stacks Project chapter file sites-cohomology.tex; no statement or proof step was rewritten, reordered or removed.\par\medskip
\noindent\textit{Editorial note.} The proof is the author's complete argument; the two displayed equalities it needs are shown, the finite case is derived from the generator property and the injectivity hypothesis, and the passage to the general presentation is carried out with exactness of filtered colimits. No step is left to the reader.\par\medskip
\noindent\textit{Editorial note.} Every labelled reference occurring inside the counted spans is carried in assembly.external\_references as a non-load-bearing navigation pointer (same-chapter ones at their definition site, cross-chapter ones at the invocation inside the counted span), so no load-bearing mathematics is replaced by a stub. No bibliography citation occurs inside the counted spans.\par\medskip
\noindent\textit{Editorial note.} The same-name check for this card scanned every staged Stacks chapter of this repository -- all sources/stacks-*/pinned-original-source.tex files, the card's own chapter excluded -- and the labels found are recorded in dedup.same\_name\_cross\_chapter on the card.\par\medskip
\noindent\textit{Editorial note.} The preamble is the chapter's own retained minimal preamble (sources/stacks-sites-cohomology/retained-preamble.tex), a verbatim excerpt of the native preamble: the Xy-pic setup and the native theorem declarations with the mathematical macros.\par\medskip
\noindent\textit{Allowed standard prerequisites (editorial).}\begin{itemize}
\item The construction of the left adjoint. Same chapter, Lemma lemma-gabriel-popescu-left-adjoint and the explicit construction of the functor F inside its proof. Supplies the description of the values of the left adjoint as a quotient of direct sums of copies of the generator, which is what the kernel computation uses.
\item The generator of the category. Same chapter, the object U and the fact that it is a generator of the Grothendieck abelian category. Supplies that a map whose composites with all maps from the generator vanish, and hence the corresponding kernel inclusion, can be tested on the generator.
\item AB5 and exact filtered colimits. Same chapter, the AB5 axiom for a Grothendieck abelian category. Supplies that filtered colimits are exact, which is what turns the verification for finite index subsets into the general statement.
\end{itemize}\medskip
\ifdefined\hypertarget\hypertarget{sourcepointer-1}{}\fi
\noindent\textit{Source reference.} Original lemma 0F5S: sources/stacks-injectives/pinned-original-source.tex, source lines 2763--2763. Navigation pointer invoked at proof line 2844; the counted span contains the pointer only and the cited result is not reproduced inside it, so the pointer is registered as a non-load-bearing navigation stub in the established form for Stacks material. No bibliography citation is involved.\par\medskip
\begin{lemma}
\label{lemma-F-G-monos}
Let $f : M \to G(A)$ be an injective map in $\text{Mod}_R$.
Then the adjoint map $f' : F(M) \to A$ is injective too.
\end{lemma}

\begin{proof}
Choose a map $R^{\oplus n} \to M$ and consider the corresponding map
$U^{\oplus n} \to F(M)$. Consider a map $v : U \to U^{\oplus n}$
such that the composition $U \to U^{\oplus n} \to F(M) \to A$ is $0$.
Then this arrow $v : U \to U^{\oplus n}$ is an element
$v$ of $R^{\oplus n}$ mapping to zero in $G(A)$. Since $f$ is injective,
we conclude that $v$ maps to zero in $M$ which means that
$U \to U^{\oplus n} \to F(M)$ is zero by construction of $F(M)$
in the proof of Lemma \ref{lemma-gabriel-popescu-left-adjoint}.
Since $U$ is a generator we conclude that
$$
\Ker(U^{\oplus n} \to F(M) \to A) = \Ker(U^{\oplus n} \to F(M))
$$
To finish the proof we choose a surjection $\bigoplus_{i \in I} R \to M$
and we consider the corresponding surjection
$$
\pi : \bigoplus\nolimits_{i \in I} U \longrightarrow F(M)
$$
To prove $f'$ is injective it suffices to show that
$\Ker(\pi) = \Ker(f' \circ \pi)$ as subobjects of $\bigoplus_{i \in I} U$.
However, now we can write $\bigoplus_{i \in I} U$ as the filtered colimit
of its subobjects $\bigoplus_{i \in I'} U$ where $I' \subset I$
ranges over the finite subsets. Since filtered colimits are
exact by AB5 for $\mathcal{A}$, we see that
$$
\Ker(\pi) =
\colim_{I' \subset I\text{ finite}}
\left(\bigoplus\nolimits_{i \in I'} U\right)
\bigcap \Ker(\pi)
$$
and
$$
\Ker(f' \circ \pi) =
\colim_{I' \subset I\text{ finite}}
\left(\bigoplus\nolimits_{i \in I'} U\right)
\bigcap \Ker(f' \circ \pi)
$$
and we get equality because the same is true for each $I'$ by
the first displayed equality above.
\end{proof}

\end{document}
